%% ---------- KODEKETA, LETRA-TIPOAK, HIZKUNTZA ----------
\usepackage[utf8]{inputenc}
\usepackage[T1]{fontenc}
\usepackage{libertine}
\usepackage{inconsolata}
\usepackage[basque]{babel}
\addto\captionsbasque{%
  \renewcommand{\contentsname}{Gaien aurkibidea}%
  \renewcommand{\listfigurename}{Irudien aurkibidea}%
  \renewcommand{\listtablename}{Taulen aurkibidea}%
  \renewcommand{\bibname}{Bibliografia}%
  \renewcommand{\figurename}{irudia}%
  \renewcommand{\tablename}{taula}%
}

%% ---------- PAKETEAK (beharrezkoak bakarrik: azkar konpilatzen du) ----------
\usepackage{microtype}
\usepackage{xcolor}
\usepackage{graphicx}
\usepackage{booktabs}          % taula txukunak
\usepackage{tabularx}          % orriaren zabalerara egokitzen diren taulak
\usepackage{amsmath}
\usepackage{tikz}              % eskemak
\usetikzlibrary{babel}         % TikZ eta babel bateragarri egiteko
\usepackage{subcaption}        % irudi bat baino gehiago batera: (a), (b)
\usepackage[font=small,labelfont=bf]{caption}
\usepackage{listings}          % kodea (minted ez bezala, ez du shell-escape behar)
\usepackage[hyperindex,bookmarks,colorlinks=true,citecolor=blue,urlcolor=blue,linkcolor=blue]{hyperref}

\graphicspath{{irudiak/}}
\setsecnumdepth{subsection}
\maxtocdepth{subsection}
\setlength{\parskip}{5pt}
\setlength{\emergencystretch}{3em}   % kode-izen luzeek lerroa ez gainditzeko

%% Oinen formatua euskaraz: "2.1. irudia", "3.1. taula", "3.1. kode-zatia"
\DeclareCaptionLabelFormat{euskaraz}{#2.~#1}
\captionsetup{labelformat=euskaraz}
\captionsetup[sub]{labelformat=parens}

%% ---------- KODE-BLOKEAK ----------
\definecolor{kodehondoa}{gray}{0.96}
\definecolor{kodeiruzkina}{rgb}{0.25,0.5,0.25}
\lstset{
  language=Python, basicstyle=\ttfamily\small, backgroundcolor=\color{kodehondoa},
  commentstyle=\color{kodeiruzkina}, keywordstyle=\color{blue!60!black},
  numbers=left, numberstyle=\tiny\color{gray}, numbersep=6pt,
  frame=single, rulecolor=\color{gray!40}, breaklines=true, tabsize=4, showstringspaces=false,
  captionpos=t, xleftmargin=1.2em, framexleftmargin=1.2em,
  inputencoding=utf8, extendedchars=true,
  literate={á}{{\'a}}1 {é}{{\'e}}1 {í}{{\'i}}1 {ó}{{\'o}}1 {ú}{{\'u}}1 {ñ}{{\~n}}1
           {Á}{{\'A}}1 {É}{{\'E}}1 {Í}{{\'I}}1 {Ó}{{\'O}}1 {Ú}{{\'U}}1 {Ñ}{{\~N}}1 {ü}{{\"u}}1
}
\renewcommand{\lstlistingname}{kode-zatia}

%% ---------- BETETZEKO LAGUNTZAK ----------
% \bete{testua}: zuek idatzi beharrekoa (grisez). Ezabatu idatzi ondoren.
\newcommand{\bete}[1]{{\color{gray}[#1]}}
% \egiaztatu{balioa}: zuen kodearekin edo datuekin bat datorren egiaztatu (laranjaz).
\definecolor{egiaztatukolorea}{rgb}{0.85,0.42,0}
\newcommand{\egiaztatu}[1]{{\color{egiaztatukolorea}#1}}
% \irudihutsa{oina}{etiketa}{zer jarri}: irudi hutsa. Ordezkatu laukia
% \includegraphics[width=.8\textwidth]{fitxategia} komandoarekin (fitxategia irudiak/ karpetan).
\newcommand{\irudihutsa}[3]{%
  \begin{figure}[htbp]\centering
    \fbox{\parbox[c][5.5cm][c]{.8\textwidth}{\centering\color{gray}[#3]}}
    \caption{#1}\label{#2}
  \end{figure}}
